\section{Finite demand closure, traces, and compatible interaction}

\subsection{Why query dependencies can be cyclic}

A finite proof may request the precision of one object as a function of the
precision available for another.  In a composition graph, local error estimates
can point in both directions.  For example, a normalization step may need a
bound on a denominator, while the denominator bound is derived from a rough
normalization.  Treating such dependencies as an acyclic syntax tree either
fails or introduces unnecessarily large global precision.

The correct finite construction is a least monotone closure.  Fix capacities
\(c_1,\ldots,c_m\) and the product of finite chains
\[
 D=\prod_{i=1}^m\{0,1,\ldots,c_i\}.
\]
Let \(d^{(0)}\) be the external demand and let \(F:D\to D\) be a monotone
transfer function collecting all local backward requirements.  Define
\[
 d^{(n+1)}=d^{(n)}\vee F(d^{(n)}),
\]
where the join is coordinatewise maximum.

\begin{theorem}[Finite monotone demand closure]
The iteration stabilizes after at most
\(\sum_i(c_i-d_i^{(0)})\) strict steps.  Its stable value \(d^*\) is the least
prefixed point above \(d^{(0)}\):
\[
 d^{(0)}\le d^*,\qquad F(d^*)\le d^*,
\]
and if \(d^{(0)}\le e\) and \(F(e)\le e\), then \(d^*\le e\).
\end{theorem}

\begin{proof}
Every iterate lies in \(D\), is above \(d^{(0)}\), and is no smaller than the
previous iterate.  If a step is strict, at least one coordinate increases by at
least one.  Therefore the rank
\[
  \rho(d)=\sum_i d_i
\]
increases at each strict step and is bounded by \(\sum_i c_i\).  This proves
stabilization within the stated budget.  At a stable value,
\(d^*=d^*\vee F(d^*)\), so \(F(d^*)\le d^*\).

For leastness, let \(e\) be any prefixed point above the seed.  Induct on
\(n\).  The base is the seed assumption.  If \(d^{(n)}\le e\), monotonicity
gives \(F(d^{(n)})\le F(e)\le e\), hence
\(d^{(n+1)}=d^{(n)}\vee F(d^{(n)})\le e\).  Thus every iterate, in particular
the stable one, lies below \(e\).
\end{proof}

For object formalization, encode the finite vector as one natural number and
make the dimension a metalevel numeral.  Coordinate projections, saturated
updates, and the rank are primitive recursive.  The invariant consists of
coordinate bounds, inflation, and the rank equation, all quantifier-free.

\subsection{Closure does not prove feasibility}

Demand closure computes the least precision request consistent with local
rules.  It does not by itself prove that every requested producer exists.  A
coordinate may saturate at its capacity while a local error target still fails.
STRIVE records this as a failed local checker or as a residual producer
formula.  This distinction is essential: a terminating scheduler is not a
mathematical existence proof.

\subsection{From closure to a clocked schedule}

Once demands are fixed, replace the cyclic graph by a finite schedule.
Each node receives its stabilized query, and a topological order is chosen for
the actual producers.  When a producer is relational rather than functional,
use a state machine whose state stores the finite information accumulated so
far.

A deterministic trace machine has primitive-recursive terms
\[
 s_0=\mathsf{init}(a),
 \qquad
 s_{t+1}=\mathsf{step}(t,s_t,a).
\]
Together with a quantifier-free invariant \(I(t,s_t,a)\).  If iPRA derives the base and
step implications, quantifier-free induction derives \(I(T,s_T,a)\) at every
finite clock \(T\).  The final checker converts the invariant into the desired
postcondition.

The same pattern works for a normalized trace certificate supplied externally.
The checker verifies the initial row, every step equation, and the terminal
condition.  Soundness follows by induction over the row index.  An external
program may generate the trace efficiently, but the proof trusts only the
checker.

\subsection{Relational transitions and backward viability}

Pointwise existence of answers is not enough to construct a compatible run.
The invalid inference is
\[
  \forall t\;\exists a_t\;\mathsf{Local}(t,a_t)
  \quad\Longrightarrow\quad
  \exists(a_0,\ldots,a_T)\;\mathsf{CompatibleRun}.
\]
Local answers may have incompatible successors.

For a bounded relational machine, let \(R_t(s,a,s')\) be the transition
checker and \(A_T(s)\) the terminal acceptance predicate.  Define viability
backwards:
\begin{align*}
 V_T(s)&=A_T(s),\\
 V_t(s)&=\mathsf{any}_{a<B_t(s)}
   \mathsf{any}_{s'<S_{t+1}}
     \bigl(R_t(s,a,s')\cdot V_{t+1}(s')\bigr).
\end{align*}
All bounds are finite primitive-recursive terms.  The value
\(V_t(s)=1\) means that a successful continuation exists within the declared
horizon and state bounds.

\begin{theorem}[Viability correctness]
For every fixed finite horizon, \(V_t(s)=1\) if and only if there is a bounded
accepted continuation from \((t,s)\).
\end{theorem}

\begin{proof}
Use backward induction on the remaining horizon.  At the terminal layer this is
the definition.  In the step, the bounded-any fold is one exactly when some
pair \((a,s')\) passes the local relation and has viable continuation.  Apply
the induction hypothesis to that continuation.  Conversely, an accepted run
supplies such a pair and hence makes the bounded-any fold one.  In a forward
\(\PRAAx\) proof, define the evaluator by primitive recursion on remaining
resource and prove the equivalence with a quantifier-free path-checker bit.
\end{proof}

A canonical compatible run is obtained by selecting the least pair whose local
edge and successor viability bits are one.  The choice is safe because it tests
future feasibility before committing to the present answer.  A failed initial
viability bit is itself a finite failure certificate: the backward table shows
that every possible continuation closes.

\subsection{Relation to signed formula certificates}

A fixed-horizon interaction is a bounded formula with alternation, so it can be
compiled directly by the signed strategy theorem.  Conversely, the formula
compiler may be implemented as a specialized backward game evaluator.  STRIVE
keeps both interfaces because their engineering uses differ:

\begin{itemize}
\item formula syntax is convenient for theorem statements and direct proof
composition;
\item a state machine is convenient for long traces, shared subcomputations,
and generated certificates;
\item viability is convenient when local witnesses must be compatible;
\item independent stable observers are cheaper when compatibility is irrelevant
to the final observation.
\end{itemize}

No logical strength is gained by choosing one representation.  The choice is a
normalization decision made before proof generation.
